% !TEX root = ../main.tex
\documentclass[12pt]{article}
\usepackage[a4paper,left=1in,right=1in,top=1in,bottom=1in]{geometry}
\usepackage{fontspec}
% 로마자는 저장소의 다른 원고처럼 Times New Roman을 쓴다. 이 글꼴이 없는 곳(Linux)에서는
% 글자 폭이 같은 Liberation Serif를 써서 줄바꿈이 같게 나오도록 한다.
\IfFontExistsTF{Times New Roman}{\setmainfont{Times New Roman}}{%
  \IfFontExistsTF{Liberation Serif}{\setmainfont{Liberation Serif}}{\setmainfont{TeX Gyre Termes}}}
% 한글 글꼴(xetexko). 본문은 명조, 산세리프는 고딕, 고정폭은 고딕코딩.
% Linux·Overleaf에는 나눔 글꼴, Windows에는 바탕·맑은 고딕을 쓴다(쉬운 한국어 번역본과 같은 방식).
\usepackage{kotex}
\IfFontExistsTF{NanumMyeongjo}{%
  \setmainhangulfont{NanumMyeongjo}[AutoFakeSlant=0.15]%
}{%
  \IfFontExistsTF{Batang}{%
    \setmainhangulfont{Batang}[AutoFakeSlant=0.15]%
  }{%
    \setmainhangulfont{Malgun Gothic}[AutoFakeSlant=0.15]%
  }%
}
\IfFontExistsTF{NanumGothic}{%
  \setsanshangulfont{NanumGothic}%
}{%
  \setsanshangulfont{Malgun Gothic}%
}
\IfFontExistsTF{NanumGothicCoding}{%
  \setmonohangulfont{NanumGothicCoding}[AutoFakeSlant=0.15]%
}{%
  \setmonohangulfont{Malgun Gothic}%
}
\usepackage{setspace}
\usepackage{amsmath}
\usepackage{amssymb}
\usepackage{amsthm}
\usepackage{booktabs}
\usepackage{array}
\usepackage{tabularx}
\usepackage{enumitem}
\usepackage{caption}
\usepackage{xcolor}
\usepackage{listings}
\usepackage{algorithm}
\usepackage{algpseudocode}
\usepackage{tikz}
\usetikzlibrary{arrows.meta,positioning}
\usepackage{hyperref}
\usepackage{xurl}
\hypersetup{
  hidelinks,
  pdftitle={Stacks 패키지 레지스트리를 위한 Clarity 스마트 컨트랙트 순위 산정: ERC-20 allowance를 대신하는 포스트컨디션},
  pdfauthor={권태홍 (Taehong Kwon)},
  pdfkeywords={Stacks, Clarity, 패키지 레지스트리, PageRank, 시빌 공격, 포스트컨디션}
}

\setstretch{1.5}
\setlength{\emergencystretch}{2em}
\setlength{\parindent}{0pt}
\setlength{\parskip}{0.55\baselineskip}
\setlist[itemize]{leftmargin=2.2em,labelsep=0.55em,itemsep=0.15\baselineskip,topsep=0.2\baselineskip,parsep=0pt}
\setlist[enumerate]{leftmargin=2.4em,labelsep=0.55em,itemsep=0.15\baselineskip,topsep=0.2\baselineskip,parsep=0pt}
\captionsetup{font=small,labelfont=bf}
% 큰 표가 따로 한 쪽을 차지하지 않고 본문 쪽에 들어가도록 한다.
\renewcommand{\topfraction}{0.85}
\renewcommand{\bottomfraction}{0.6}
\renewcommand{\textfraction}{0.12}
\renewcommand{\floatpagefraction}{0.7}
\renewcommand{\arraystretch}{1.15}

% 한국어 이름표: 그림, 표, 코드, 참고문헌, 부록, 증명, 알고리즘
\AtBeginDocument{%
  \renewcommand{\figurename}{그림}%
  \renewcommand{\tablename}{표}%
  \renewcommand{\lstlistingname}{코드}%
  \renewcommand{\refname}{참고문헌}%
  \renewcommand{\appendixname}{부록}%
  \renewcommand{\proofname}{증명}%
}
\floatname{algorithm}{알고리즘}
\algrenewcommand\algorithmicrequire{\textbf{입력:}}
\algrenewcommand\algorithmicensure{\textbf{출력:}}
\algrenewcommand\algorithmicreturn{\textbf{반환}}
\algrenewcommand\algorithmicrepeat{\textbf{반복}}
\algrenewcommand\algorithmicuntil{\textbf{멈춤 조건:}}
\algrenewtext{For}[1]{\textbf{각각에 대해:} #1}
\algrenewtext{EndFor}{\textbf{끝}}

\newtheorem{proposition}{명제}
\theoremstyle{definition}
\newtheorem{definition}{정의}

% 본문 전체에서 쓰는 이름
\newcommand{\PCER}{PC-EndorseRank}
\newcommand{\cpm}{\texttt{cpm}}
% 본문 안의 코드. Clarity 이름은 길고 하이픈이 많으므로 하이픈, 점, 밑줄 뒤에서
% 줄을 바꿀 수 있게 한다(줄이 바뀌어도 하이픈을 덧붙이지 않는다).
\ExplSyntaxOn
\tl_new:N \l__ms_code_tl
\NewDocumentCommand{\code}{m}
 {
  \tl_set:Nn \l__ms_code_tl {#1}
  \tl_replace_all:Nnn \l__ms_code_tl {-} {-\allowbreak}
  \tl_replace_all:Nnn \l__ms_code_tl {.} {.\allowbreak}
  \tl_replace_all:Nnn \l__ms_code_tl {\_} {\_\allowbreak}
  \texttt{\tl_use:N \l__ms_code_tl}
 }
\ExplSyntaxOff
\newcommand{\R}{\mathbb{R}}

% Clarity 코드 목록: 고정폭 글꼴, 키워드는 굵게
\lstdefinelanguage{Clarity}{
  sensitive=true,
  alsoletter={-?!},
  morekeywords={define-public,define-private,define-read-only,define-constant,define-map,
    define-data-var,define-trait,impl-trait,use-trait,contract-call?,as-contract?,
    restrict-assets?,with-ft,with-stx,with-nft,with-stacking,with-staking,with-pox,with-all-assets-unsafe,
    try!,unwrap!,asserts!,begin,let,ok,err,tx-sender,contract-caller,current-contract},
  morecomment=[l]{;;},
  morestring=[b]",
}
\lstset{
  basicstyle=\ttfamily\small,
  keywordstyle=\bfseries,
  commentstyle=\itshape,
  columns=fullflexible,
  keepspaces=true,
  breaklines=true,
  frame=single,
  framesep=4pt,
  xleftmargin=6pt,
  xrightmargin=6pt,
  aboveskip=0.6\baselineskip,
  belowskip=0.4\baselineskip,
  showstringspaces=false,
}
